% ─── EXAMPLE DRAFT: Related Work section ───────────────────────────────────
% Reference draft for structure/flow — NOT meant to be copied verbatim.
% Target length: ~1–1.5 columns IEEE two-column format, + table*.
% See README.md in this directory for full context and open decisions.
%
% STATUS: PyJama dropped (see README). Hameed/Ziemann inclusion still open.
% ───────────────────────────────────────────────────────────────────────────

\section{Related Work}

Recent surveys~\cite{jamming_survey_2024} document the shift from fixed-heuristic
jamming toward learning-based strategies. We organize prior work along the three
axes central to our contribution: cooperative multi-agent jamming,
deep-learning-based jamming detection, and generative models for adversarial
signal synthesis.

\paragraph{Cooperative MARL for jamming}
Reinforcement learning has been applied to both attack and defense in wireless
systems, but predominantly with discrete or low-dimensional action spaces.
Zhang and Wu~\cite{article} coordinate multiple jammers against a
frequency-hopping network via hierarchical DRL modeled as a Stackelberg game,
achieving higher success rates than non-cooperative baselines; however, each
jammer selects from a finite set of frequency--power combinations rather than
synthesizing waveforms.
Valianti et al.~\cite{valianti2024cooperative} employ cooperative Q-learning to
neutralize rogue drone swarms with discrete power levels.
On the defense side, Abolhassani et al.~\cite{abolhassani2025coordinated} use
QMIX for coordinated anti-jamming spectrum access, training against simplistic
energy-threshold adversaries that ignore the structure of transmitted data.
The CTDE paradigm introduced by Lowe et al.~\cite{NIPS2017_68a97503} and the
MAPPO algorithm of Yu et al.~\cite{yu2022mappo} provide well-validated
multi-agent training foundations, yet no prior work applies cooperative MARL to
continuous IQ-level waveform synthesis for jamming.

\paragraph{Deep-learning-based jamming detection}
CNN classifiers on spectrograms or raw IQ samples represent the current
state of the art in jamming detection.
Li et al.~\cite{li2022jamming_ofdm} achieve 99.8\% accuracy on OFDM
spectrograms using EfficientNet-B0, a result we replicate as our frozen
adversarial detector.
Erpek et al.~\cite{erpek2019deep} frame jamming as an adversarial ML problem,
training CNNs to both detect attacks and improve jammer strategies.
% OPEN: include Hameed et al. here? Transmitter-side evasion, not jamming,
% but directly demonstrates "shape waveform statistics to evade a learned
% classifier" — closest conceptual match to the stealth objective.
% In parallel, Hameed et al.~\cite{hameed2021offense} show that a legitimate
% transmitter can perturb its own symbols to evade a DL-based modulation
% classifier while preserving decodability---demonstrating that adversarial
% waveform shaping can defeat learned detectors.
These results motivate a jammer that shapes its signal statistics to evade
detection, a capability absent from all cooperative MARL jamming work above.

\paragraph{Generative models for adversarial wireless signals}
Shi et al.~\cite{sagduyu2021gan} use a GAN to synthesize spoofing IQ signals
over the air that fool a deep-learning classifier, demonstrating that generative
models can produce realistic adversarial wireless waveforms.
Wen et al.~\cite{wen2025generative} extend this idea to cooperative jammers with
GAN-aided covert communication.
However, GANs cannot provide the exact action log-probabilities required by
policy-gradient methods such as PPO, precluding their use as MARL policy heads.
% OPEN: include Ziemann & Metzler here? Radar domain (LPD waveform design),
% not communications, but it's the strongest existing validation of the
% "generative model for dual-objective (effective + undetectable) waveform
% design" paradigm you're applying to comms. Cross-domain cite — your call
% on whether it strengthens or distracts.
% In the radar domain, Ziemann and Metzler~\cite{ziemann2024lpd} show that a
% generative model can synthesize waveforms that are simultaneously effective
% for sensing and statistically indistinguishable from the RF background,
% validating the dual-objective (effective + undetectable) waveform design
% paradigm we adopt for communications.
Ward et al.~\cite{ward2019nf_rl} demonstrate that normalizing-flow policies
capture multimodal action distributions and outperform Gaussian baselines in
continuous control, providing the policy architecture we extend to the
multi-agent wireless domain.

\paragraph{Research gap}
As summarized in Table~\ref{tab:related_works}, existing work addresses
individual aspects of intelligent jamming but no prior method combines all three:
(i)~IQ-level waveform generation via expressive generative models,
(ii)~cooperative multi-agent coordination, and
(iii)~learned stealth against a neural-network detector.
We fill this gap by equipping MAPPO agents with neural spline flow
(NSF)~\cite{durkan2019nsf} policies that synthesize complex-valued OFDM jamming
signals, jointly optimizing for disruption and detection avoidance under a
black-box threat model.

% ─── Comparison table ──────────────────────────────────────────────────────
\begin{table*}[!t]
\renewcommand{\arraystretch}{1.3}
\caption{Comparison of related work along the three novelty axes of the proposed
method. ``IQ Waveform'' indicates whether the method generates raw complex-valued
signals rather than selecting from discrete actions.}
\label{tab:related_works}
\centering
\begin{tabular}{l l c c c c}
\hline\hline
\textbf{Reference} & \textbf{Method} & \textbf{IQ Waveform} & \textbf{MARL}
  & \textbf{Cooperative} & \textbf{Det.\ Avoidance} \\
\hline
Zhang \& Wu~\cite{article}
  & HDRL (Stackelberg)     & No  & Yes & Yes & No  \\
Valianti \emph{et al.}~\cite{valianti2024cooperative}
  & Cooperative Q-learning & No  & Yes & Yes & No  \\
Abolhassani \emph{et al.}~\cite{abolhassani2025coordinated}
  & QMIX (defense)         & No  & Yes & Yes & No  \\
Shi \emph{et al.}~\cite{sagduyu2021gan}
  & GAN spoofing           & Yes & No  & No  & Partial \\
Wen \emph{et al.}~\cite{wen2025generative}
  & GAN-aided covert comm. & No  & No  & Yes & Yes \\
% Ziemann \& Metzler~\cite{ziemann2024lpd} & GAN (radar LPD) & Yes & No & No & Yes \\  % OPEN: include?
\textbf{Proposed}
  & \textbf{MAPPO + NSF}   & \textbf{Yes} & \textbf{Yes}
  & \textbf{Yes} & \textbf{Yes} \\
\hline\hline
\end{tabular}
\end{table*}
